\section{The shape of the proof}\label{sec:overview}

The top of the proof is short. \Cref{st:rh} follows from
\cref{st:grh} for the trivial character by the functional equation, in a few dozen lines on
Mathlib's interface. \Cref{st:grh} is obtained for every character from a statement
\lean{ZeroFreeAll} about primitive characters. A theorem
(\lean{zero_free_all_of_green_v_win}) derives that statement from three hypotheses, and with
the passage to all characters it reads
\begin{gather*}
\text{\lean{every_character_of_green_v_win}}\colon\\
\mathsf{Green}\ \longrightarrow\ \mathsf{VWin}\ \longrightarrow\ \mathsf{Nonempty}(\mathsf{HypsAlt})
\ \longrightarrow\ \mathsf{EveryCharacter},
\end{gather*}
each of which is in turn proved in the archive. They are the three inputs of the argument:
\begin{itemize}
\item $\mathsf{Green}$: a statement about the characters of finite general linear groups, in
  the parametrisation of J.~A.~Green \cite{green1955}, used to count local types;
\item $\mathsf{VWin}$: a bound on Galois cohomology groups attached to an auxiliary field,
  valid on windows of scales (\cref{sec:cohomology});
\item $\mathsf{HypsAlt}$: a record of forty-six fields about families of lattices of
  automorphic forms on unitary groups (\cref{sec:lattices,sec:frame,sec:literature}).
\end{itemize}

The body of that theorem \lvP{} applies a windowed form of Landau's theorem, on windows
$[Y,Y^{3/2}]$, to a statement $\mathsf{TopG}$, which says: for every small
$\eps>0$ there is a family of arithmetic data, indexed by scales $x$ in long windows, on which
a condition \lean{ChainTied} holds. \lean{ChainTied} packages eleven inequalities between real
numbers (\cref{sec:lattices}), with size conditions and with the identification of those
numbers with the arithmetic data of the count and with slopes and heights of lattices. The
inequalities, with the size conditions and the identification of the count, imply
$G(x)\le x^{1/2+C\eps}$ by an argument that involves
no arithmetic beyond an identity of character sums (\cref{sec:analytic}).

So the proof has two halves that meet in eleven real inequalities:
\begin{center}
\begin{tabular}{@{}p{0.46\textwidth}@{\qquad}p{0.46\textwidth}@{}}
\toprule
\emph{Analytic half} (\cref{sec:analytic}) & \emph{Arithmetic half} (Sections~\ref{sec:lattices}--\ref{sec:hosts})\\
\midrule
Landau on windows; an exact identity for a count of prime powers; averaging; balancing three
error terms at $m\approx\sqrt x$.
&
Lattices of forms, period maps, divisibility at inert primes, a slope inequality, a period
computation, a trace formula, a cohomological bound, auxiliary fields.\\
\bottomrule
\end{tabular}
\end{center}
The analytic half is elementary. Almost all of the archive is the
arithmetic half and the published theorems it formalises.

\subsection*{Notation}
In the account of the argument (\cref{sec:analytic,sec:lattices,sec:divisibility,sec:frame})
the letters of the following table have the meaning given here, except where one is the index
or the bound variable of a formula, or carries a subscript other than those of the table.
\Cref{sec:cohomology,sec:literature} describe inputs in the
notation of their sources, and there $E$, $G$, $H$, $\lambda$, $h$, $n$, $o$, $p$ and $q$ are
also used as the sources use them. Any other letter is local to the statement or paragraph in
which it occurs.
\begin{center}\small
\begin{tabular}{@{}p{0.2\textwidth}p{0.74\textwidth}@{}}
\toprule
$x$, $\eps$, $b$ & the scale, a large real number; a small fixed exponent; $b=\exp(x^\eps)$\\
$m$, $k$ & the number of lifts, $x^{1/2-\eps}\le m\le x^{1/2+\eps/2}$; the rank parameter,
  $xm\le k\le n^2x^{1/2-\eps}$\\
$E$, $\mathfrak G$, $H$, $n$ & the host field; its Galois group $(\Z/2)^2\times H$ over $\Q$;
  $n=\abs H$\\
$\Phi$, $\chi_\Phi$, $f$ & the CM type, a subgroup of index two of $\mathfrak G$; its quadratic
  character; $f=\abs\Phi=2n$, the number of infinite places of $E$\\
$q$, $d$, $e_2$, $\sigma$ & the two primes of the host and the exponent of its $2$-power layer
  (\cref{sec:hosts}); the non-trivial element of the inertia group at $d$\\
$G_\chi$, $G$ & the smoothed sum over prime powers attached to a character
  \eqref{eq:G}; $G=G_1$\\
$\mathcal L$, $\mathcal K$, $\psi_{\mathcal K}$, $o$ & a level: its set of characters, its
  distinguished character, its integer \eqref{eq:level}\\
$A_0$, $TT$, $T_2$, $\delta_T$ & a subgroup of $H$; the set of top places; the thick set;
  its density $\abs{T_2}/n$\\
$t$, $a_t$, $w$, $u_0$ & a top place, an element of $TT$; its size; a system of weights on
  $\Phi$; a translate in $A_0$\\
$\mathcal K'$, $\beta$ & the kept characters and their common weight (\cref{sec:reduction})\\
$\ASK$, $\MAIN$, $\Err$ & the count \eqref{eq:ask}, its main term \eqref{eq:main}, the error
  of \cref{sec:balance}\\
$u$, $M_u$, $N_u$, $B_u$ & a word; its lattice of forms; the rescaled sublattice; its period
  map (\cref{sec:cast})\\
$h(u)$, $T(u)$, $P$ & the height bound of $B_u$; the level prime; the lattice of periods\\
$\lambda$, $g(u,\lambda)$ & a finite place of the coefficient field; the exponent of the
  rescaling there, a sum of gains and corrections \eqref{eq:Nu}\\
$p$, $h_p$, $\gamma_p$ & a rational prime; the component in $H$ of its Frobenius element;
  $\gamma_p=p/(p^2-1)$\\
$\ell$ & the prime of \cref{sec:divisibility} and of \cref{sec:cohomology}\\
$\gamma$, $\Omega$, $\ph$ & Euler's constant; a CM period; $\ph(y)=y\log y-y$, so that
  $\ph(a)=\log a!+O(\log a)$\\
$\bar\mu$, $HH$, $u_F$, $\delta_F$, $\eta_F$, $e_F$ & the parameters of the slope inequality
  (\cref{st:F})\\
$t_0$, $\len_\ell(e)$ & the number of twists $\lfloor x^{401}\rfloor$; the length of
  $\Sha(E,\ell,e)$ (\cref{sec:cohomology})\\
\bottomrule
\end{tabular}
\end{center}
The weights of a frame are written $\gamma_i$ with an index.
